% Chapter 6: the minimum-amount guarantee.
\chapter{The minimum-amount guarantee}\label{ch:guarantee}
\index{guarantee}\index{margin}

The machine is built so that a customer never gets less paper than the label
says. Every wind ends in exactly one of two ways:

\begin{center}
\renewcommand{\arraystretch}{1.5}
\begin{tabularx}{\textwidth}{@{}>{\raggedright\arraybackslash}p{0.28\textwidth}>{\raggedright\arraybackslash}X@{}}
\toprule
\textbf{Result} & \textbf{What it promises} \\
\midrule
\shows{DONE}, three short beeps &
  At least the full length \textbf{plus the margin} (1 inch from the factory)
  went past the measuring wheel. The roll is good. \\
\shows{!! UNDER TARGET !!} &
  The machine could \textbf{not} confirm the full length. Treat the roll as
  short. Alarm beeps unless you stopped the wind yourself. \\
\bottomrule
\end{tabularx}
\end{center}

\section{How it works}
\begin{enumerate}
  \item The \textbf{goal} is the target length plus the margin. For a 30 yard
        roll that is 30 yards and 1 inch.
  \item The motor runs until the paper is one \emph{coast} (12 inches) short
        of the goal, then it is switched off and the roll glides on.
  \item After the paper has stopped, the machine compares the count with the
        goal. If it is short, it bumps the motor briefly and checks again.
  \item It only says \shows{DONE} when the stopped count is at or past the
        goal. It never guesses: the decision is made on paper that has stopped
        moving.
\end{enumerate}

\section{What the guarantee depends on}
The guarantee is only as good as the measurement. It assumes:
\begin{itemize}
  \item The measuring wheel is \textbf{calibrated} (Chapter~\ref{ch:calibration})
        and its rubber is clean and not worn.
  \item The wheel turns with the paper \textbf{without slipping} and is not
        lifted off the paper.
  \item The paper goes onto the roll; a tear or a jam upstream of the wheel
        will not be noticed while the motor is running (only during the trim).
\end{itemize}

\section{How much over?}
A normal roll is a little longer than the label, usually by a few inches,
because the roll coasts on after the last bump. The design estimate is about
1.3 inches over on average and no more than about 14 inches in the worst
case. These figures are estimates from the design and have not yet been
measured on the bench (Chapter~\ref{ch:limits}). If your rolls are regularly
much longer than that, recalibrate the coast (Chapter~\ref{ch:calibration}).

\begin{note}[Tip]
The \shows{DONE} screen tells you the exact overrun of every roll, for
example \shows{+4.3 in over}. Glance at it each time: a sudden change is the
first sign of a dirty or slipping wheel.
\end{note}
